\reportheading{4.\quad Task 2 --- Smoke extraction design and optimisation}
\subhead{4.1\quad FDS screening of the original 16 designs}
The best of the 16 original OP-6 designs combined a 0.16~m curtain with 0.32~m$^3$/s extraction at baseline soot yield (Appendix~C). Even this design gave a minimum interface of 0.095~m, leaving a deficit of 0.285~m against the 0.38~m target. The next trials therefore changed where replacement air entered and where the ceiling vents collected smoke.

\subhead{4.2\quad Locating the failure and adapting the inlets}
An OP-6 trial used 0.36~m end inlets, 0.24~m curtains and 1.28~m$^3$/s total extraction. At soot yield 0.05 the far-end column ($x=6.1$~m) recorded an interface of 0.353~m at 83.0~s. The readings immediately before and after both exceeded 0.60~m. The low reading counts as a breach of the thermal criterion. At the same point visibility at half height remained 30~m and the vertical temperature gradient was weak, leaving the extent of smoke there uncertain.

The choice of inlet height for OP-6 can be seen by comparing the two inlet heights at each of two soot yields on the original grid. At 1.28~m$^3$/s, raising the end openings from 0.44 to 0.48~m increased the minimum interface from 0.431 to 0.583~m at yield 0.02, but reduced it from 0.524 to 0.497~m at yield 0.05. Taking the lower reading for each inlet height, the taller opening raised the limiting value from 0.431 to 0.497~m. We therefore used 0.48~m for OP-6 in the six-scenario check, with 0.44~m for OP-1 and 0.24~m curtains in both states.
\reportfigure{\begin{tikzpicture}
\begin{axis}[width=.94\linewidth,height=5.6cm,xmin=.009,xmax=.112,ymin=.22,ymax=.405,
xtick={.02,.05,.10},ytick={.24,.28,.32,.36,.38,.40},scaled ticks=false,
xticklabel style={/pgf/number format/fixed,/pgf/number format/fixed zerofill,/pgf/number format/precision=2},
yticklabel style={/pgf/number format/fixed},axis lines*=left,ymajorgrids,grid style={ink!13},
xlabel={Soot yield (kg/kg)},ylabel={Minimum height (m)},legend columns=3,legend style={at={(.5,1.03)},anchor=south,draw=none,/tikz/every even column/.append style={column sep=6pt}}]
\addplot[only marks,teal,mark=*,mark size=3pt] coordinates {(0.02,0.2627491) (0.05,0.37114426) (0.1,0.24450079)};\addlegendentry{0.16 m curtain}
\addplot[only marks,rust,mark=triangle*,mark size=3pt] coordinates {(0.02,0.36496348) (0.05,0.3527722) (0.1,0.36869806)};\addlegendentry{0.24 m curtain}
\addplot[only marks,ink,mark=square*,mark size=3pt] coordinates {(0.10,0.26696128)};\addlegendentry{0.32 m curtain}
\addplot[ink!75,dashed,forget plot] coordinates {(.009,.38) (.112,.38)};
\node[anchor=south west,font=\sffamily\footnotesize,text=ink] at (axis cs:.011,.38){0.38 m target};
\end{axis}\end{tikzpicture}}{Figure 3. Minimum thermal layer height for the seven completed OP-6 curtain trials at 1.28 m$^3$/s. Each marker represents eight stations over 120 s. The 0.32 m curtain was evaluated at soot yield 0.10 kg/kg.}

\clearpage
\subhead{4.3\quad Selecting the design flow}
With two ceiling vents, 0.24~m curtains and the inlet height set for each door state, the remaining choice was the common extraction rate. We tested 1.12~m$^3$/s, halfway between the original settings of 0.96 and 1.28~m$^3$/s, to see whether it could maintain the interface above half height under both door states at all three soot yields.
\reportfigure{\begin{tikzpicture}
\pgfplotsset{flowpanel/.style={
 scale only axis,height=3.35cm,ymin=.35,ymax=.62,ytick={.38,.44,.50,.56,.62},
 scaled ticks=false,yticklabel style={/pgf/number format/fixed,/pgf/number format/precision=2},
 ymajorgrids,grid style={ink!13},tick align=outside,title style={font=\sffamily\small\bfseries},
 tick label style={font=\sffamily\scriptsize},label style={font=\sffamily\scriptsize}}}
\begin{axis}[flowpanel,name=original,width=.47\linewidth,xmin=.55,xmax=6.45,
 title={(a) Original grid},ylabel={Minimum height (m)},
 xtick={1,2,3,4,5,6},xticklabels={0.02,0.05,0.10,0.02,0.05,0.10},
 clip=false,legend style={at={(.50,.99)},anchor=north,font=\sffamily\scriptsize,fill=white},
 legend columns=2]
\addplot[only marks,mark=*,teal,mark size=1.9pt] coordinates {(1,.45742963) (2,.47169434) (3,.44104958) (4,.39573017) (5,.52617570) (6,.40609647)};
\addlegendentry{0.96}
\addplot[only marks,mark=triangle*,rust,mark size=2.4pt] coordinates {(1,.44252295) (2,.50930185) (3,.50127480) (4,.58267633) (5,.49747471) (6,.55033156)};
\addlegendentry{1.28}
\addplot[ink!70,dashed] coordinates {(.55,.38)(6.45,.38)};
\node[anchor=north,font=\sffamily\scriptsize] at (axis description cs:.25,-.15) {OP-1};
\node[anchor=north,font=\sffamily\scriptsize] at (axis description cs:.75,-.15) {OP-6};
\end{axis}
\begin{axis}[flowpanel,at={(original.east)},anchor=west,xshift=.065\linewidth,
 width=.34\linewidth,xmin=.65,xmax=3.35,title={(b) Refined grid, OP-6},
 xtick={1,2,3},xticklabels={0.96,1.12,1.28},xlabel={Flow (m$^3$/s)},yticklabels=\empty]
\addplot[ink!70,dashed] coordinates {(.65,.38)(3.35,.38)};
\addplot[only marks,mark=*,teal,mark size=2.3pt] coordinates {(1,.37169)};
\addplot[only marks,mark=diamond*,ink,mark size=3pt] coordinates {(2,.5522)};
\addplot[only marks,mark=triangle*,rust,mark size=2.6pt] coordinates {(3,.5429)};
\node[anchor=south,font=\sffamily\scriptsize] at (axis cs:2,.558) {0.5522 m};
\end{axis}
\end{tikzpicture}}{Figure 4. The dashed line marks the 0.38~m thermal-interface target. (a) Two reference flows (m$^3$/s) on the original grid, grouped by door state and soot yield (kg/kg). (b) Three flows on the same refined grid for OP-6 at yield 0.02. The 1.12 and 1.28~m$^3$/s minima cover 120~s. The 0.96~m$^3$/s record ends at 98~s. Figure~5 gives the six-scenario check at 1.12~m$^3$/s.}


Baseline soot gave the smallest original-grid margin at both reference flows (Figure~4a). At 0.96~m$^3$/s, only 15.7~mm separated the lowest interface from the target, within one 40~mm vertical cell. When the cell height was halved, the far-end column in the OP-6 baseline case reached 0.3717~m at 98.0~s and screening stopped (Table~D2). Figure~4b compares the three flows at this finer spacing.

\subhead{Clearance and power at 1.12~m$^3$/s}
For OP-6 at baseline soot yield, the lowest interface over 120~s was 0.5522~m at 1.12~m$^3$/s and 0.5429~m at 1.28~m$^3$/s. At the proposed setting each branch carries 0.56~m$^3$/s, giving a velocity of 1.75~m/s and an estimated pressure drop of 3.70~Pa. At the assumed fan efficiency of 0.60, the two branches together require 6.91~W. This is about 33\% below the 10.32~W at 1.28~m$^3$/s, so we carried 1.12~m$^3$/s into the six-scenario check.

\clearpage
\subhead{Six-scenario FDS confirmation}
In all six refined-grid runs at 1.12~m$^3$/s, the interface stayed above 0.38~m at each of the eight stations for the full 120~s (Table~D3). Mean heat release over 90--120~s was 102.1--103.9~kW against the nominal 106.2~kW. The far-end column governed clearance. In OP-6 at soot yield 0.05, it recorded 0.4358~m at $x=6.1$~m and 113.0~s, leaving 55.8~mm above the target. There the longitudinal velocity changed sign between $z=0.42$ and 0.44~m. Figure~A2 shows smoke at half height at the same time.
% Rows are the six refined-grid FDS station minima at 1.12 m3/s.
% PINN points are the forecasts frozen in hybrid_frozen_1p12.json.
\reportfigure{\begin{tikzpicture}
\begin{axis}[width=.96\linewidth,height=5.1cm,xmin=.55,xmax=6.45,
 ymin=.355,ymax=.61,xtick={1,2,3,4,5,6},
 xticklabels={{OP-1\\0.02},{OP-1\\0.05},{OP-1\\0.10},{OP-6\\0.02},{OP-6\\0.05},{OP-6\\0.10}},
 xticklabel style={align=center},ytick={.38,.44,.50,.56,.60},ylabel={Minimum height (m)},
 ymajorgrids,unbounded coords=discard,
 legend style={at={(.5,1.03)},anchor=south,legend columns=3,draw=none},
 /tikz/every even column/.append style={column sep=5pt}]
\addplot[only marks,mark=*,teal,mark size=2.7pt]
 table[col sep=comma,x=position,y=fds_m]{data/verified-flow-forecast.csv};
\addlegendentry{Refined-grid FDS}
\addplot[only marks,mark=diamond*,rust,mark options={fill=white},mark size=3.2pt]
 table[col sep=comma,x expr=\thisrow{position}+.10,y=pinn_m]{data/verified-flow-forecast.csv};
\addlegendentry{Frozen PINN}
\addplot[ink,densely dashed] coordinates {(.55,.38)(6.45,.38)};
\addlegendentry{Height target}
\draw[teal,thin] (axis cs:5,.385)--(axis cs:5,.430);
\node[anchor=west,font=\sffamily\scriptsize,text=teal] at (axis cs:5.04,.408) {55.8 mm};
\end{axis}\end{tikzpicture}}{Figure 5. Refined-grid verification at 1.12~m$^3$/s by door state and soot yield (kg/kg). Filled circles are the FDS minima over eight stations and 120~s. Open diamonds show the three PINN forecasts frozen before the corresponding runs.}


\subhead{4.4\quad CO generation and local smoke conditions}
Each kilogram of fuel consumed produces 10~g of CO at the specified yield. In earlier OP-6 trials with 0.24~m curtains and 1.28~m$^3$/s extraction, integrated fire power gave estimated CO masses of 2.76, 2.81 and 2.94~g as soot yield rose. Peak half-height concentrations were 20.0, 26.2 and 163.5~ppm. Capture near the plume and a clear make-up air path help remove CO. In the original OP-6 layout with baseline soot and 0.16~m curtains, increasing extraction from 0.16 to 0.64~m$^3$/s reduced mean far-end CO from 101.04 to 0.62~ppm. Table~1 gives local exposure for the recommended layout.

\tabletitle{Table 1. Half-height exposure indicators at the selected 1.12 m$^3$/s setting. Each value is the extreme over eight probes and 120 s on the refined grid.}
\begin{reporttable}{L{.18}L{.24}L{.31}L{.27}}
\rowcolor{tablegrey}Door&Soot (kg/kg)&Min. visibility (m)&Max. CO (ppm)\\
OP-1&0.02&30.00&0.65\\
OP-1&0.05&30.00&0.26\\
OP-1&0.10&30.00&0.28\\
OP-6&0.02&24.15&7.59\\
OP-6&0.05&16.42&11.81\\
OP-6&0.10&4.06&32.83\\
\end{reporttable}
In OP-6 at yield 0.10, the probe near the burner ($x=3.6$~m) recorded both peak CO and minimum visibility at 3.0~s. The thermal minimum occurred much later and farther away, at $x=6.1$~m and 113.0~s in OP-6 at yield 0.05.
